%  LaTeX support: latex@mdpi.com 
%  For support, please attach all files needed for compiling as well as the log file, and specify your operating system, LaTeX version, and LaTeX editor.

%=================================================================
\documentclass[informatics,article,submit,moreauthors,pdftex]{Definitions/mdpi} 

%--------------------
% Class Options:
%--------------------

%---------
% article
%---------
% The default type of manuscript is "article", but can be replaced by: 
% abstract, addendum, article, book, bookreview, briefreport, casereport, comment, commentary, communication, conferenceproceedings, correction, conferencereport, entry, expressionofconcern, extendedabstract, datadescriptor, editorial, essay, erratum, hypothesis, interestingimage, obituary, opinion, projectreport, reply, retraction, review, perspective, protocol, shortnote, studyprotocol, systematicreview, supfile, technicalnote, viewpoint, guidelines, registeredreport, tutorial
% supfile = supplementary materials

%----------
% submit
%----------
% The class option "submit" will be changed to "accept" by the Editorial Office when the paper is accepted. This will only make changes to the frontpage (e.g., the logo of the journal will get visible), the headings, and the copyright information. Also, line numbering will be removed. Journal info and pagination for accepted papers will also be assigned by the Editorial Office.

%---------
% pdftex
%---------
% The option pdftex is for use with pdfLaTeX. If eps figures are used, remove the option pdftex and use LaTeX and dvi2pdf.

%=================================================================
% MDPI internal commands
\firstpage{1} 
\makeatletter 
\setcounter{page}{\@firstpage} 
\makeatother
\pubvolume{1}
\issuenum{1}
\articlenumber{0}
%\doinum{}
\pubyear{2021}
\copyrightyear{2020}
%\externaleditor{Academic Editor: Firstname Lastname} % For journal Automation, please change Academic Editor to "Communicated by"
\datereceived{} 
\dateaccepted{} 
\datepublished{} 
\hreflink{https://doi.org/} % If needed use \linebreak

%====================== my packages: start ===============================
%\bibliographystyle{numeric}
\usepackage{lipsum} 
\bibstyle{numeric}
\usepackage[super]{nth}
\usepackage{subcaption}
\usepackage{float}
\usepackage{wrapfig}
\usepackage{amsmath}
\usepackage{hyperref}
\graphicspath{{figures/}} % path to folder with figures

\usepackage{acro} %list of abbreviations
\acsetup{first-style=long-short,list/display = used}
\DeclareAcronym{GPS}{ 
    short = {GPS}, 
    long  = {Global Positioning System},
    tag = {abbrev}
}
\DeclareAcronym{AI}{
    short = {AI},
    long  = {Artificial Intelligence},
    tag = {abbrev}
}
\DeclareAcronym{WWSSN}{ 
    short = {WWSSN}, 
    long  = {World-Wide Standardized Seismograph Network},
    tag = {abbrev}
}
\DeclareAcronym{EO}{
    short = {EO},
    long  = {Earth Observation},
    tag = {abbrev}
}
\DeclareAcronym{UCC}{
    short = {UCC},
    long  = {Uccle seismic station},
    tag = {abbrev}
}
\DeclareAcronym{ML}{
    short = {ML},
    long  = {Machine Learning},
    tag = {abbrev}
}
\DeclareAcronym{DL}{
    short = {DL},
    long  = {Deep Learning},
    tag = {abbrev}
}
\DeclareAcronym{FFT}{
    short = {FFT},
    long  = {Fast Fourier Transform},
    tag = {abbrev}
}
\DeclareAcronym{DFT}{
    short = {DFT},
    long  = {Discrete Fourier Transform},
    tag = {abbrev}
}
\DeclareAcronym{IFT}{
    short = {IFT},
    long  = {Inverse Fourier Transform},
    tag = {abbrev}
}
\DeclareAcronym{HSV}{
    short = {HSV},
    long  = {Hue, Saturation, Value},
    tag = {abbrev}
}







 %define abbreviations in this file.

%====================== my packages:
%%% FORMATTING STUFF
% JUST FOR SUBMITTING:
\usepackage{xcolor}
\newcommand{\review}[1]{\textcolor{blue}{#1}}
\newcommand{\todo}[1]{\textcolor{purple}{#1}}
\newcommand{\CITE}[1]{\textcolor{red}{(CITE: #1)}}

%------------------------------------------------------------------
% The following line should be uncommented if the LaTeX file is uploaded to arXiv.org
%\pdfoutput=1

%=================================================================
% Add packages and commands here. The following packages are loaded in our class file: fontenc, inputenc, calc, indentfirst, fancyhdr, graphicx, epstopdf, lastpage, ifthen, lineno, float, amsmath, setspace, enumitem, mathpazo, booktabs, titlesec, etoolbox, tabto, xcolor, soul, multirow, microtype, tikz, totcount, changepage, paracol, attrib, upgreek, cleveref, amsthm, hyphenat, natbib, hyperref, footmisc, url, geometry, newfloat, caption

%=================================================================
%% Please use the following mathematics environments: Theorem, Lemma, Corollary, Proposition, Characterization, Property, Problem, Example, ExamplesandDefinitions, Hypothesis, Remark, Definition, Notation, Assumption
%% For proofs, please use the proof environment (the amsthm package is loaded by the MDPI class).

%=================================================================
% Full title of the paper (Capitalized)
\Title{Evaluation of DigitSeis Software for Building a Vectorised Dataset of Historical Seismograms}

% MDPI internal command: Title for citation in the left column
\TitleCitation{Evaluation of DigitSeis Software for Building a Vectorised Dataset of Historical Seismograms}

% Author Orchid ID: enter ID or remove command
\newcommand{\orcidauthorA}{0000-0000-0000-000X} % Add \orcidA{} behind the author's name
%\newcommand{\orcidauthorB}{0000-0000-0000-000X} % Add \orcidB{} behind the author's name

% Authors, for the paper (add full first names)
\Author{Poina Lemenkova $^{1}$* and Olivier Debeir $^{1}$}

% MDPI internal command: Authors, for metadata in PDF
\AuthorNames{Poina Lemenkova and Olivier Debeir}

% MDPI internal command: Authors, for citation in the left column
\AuthorCitation{Lemenkova, P.; Debeir, O.}
% If this is a Chicago style journal: Lastname, Firstname, Firstname Lastname, and Firstname Lastname.

% Affiliations / Addresses (Add [1] after \address if there is only one affiliation.)
\address{%
$^{1}$ \quad Université Libre de Bruxelles, École polytechnique de Bruxelles (Brussels Faculty of Engineering), Laboratory of Image Synthesis and Analysis (LISA). Building L, Campus de Solbosch, Avenue Franklin Roosevelt 50, Brussels 1000, Belgium. }

% Contact information of the corresponding author
\corres{Correspondence: polina.lemenkova@ulb.be; Tel.: +32 471 86 04 59 (P.L.)}

% Current address and/or shared authorship
% \firstnote{Current address: Affiliation 3} 
%\secondnote{These authors contributed equally to this work.}
% The commands \thirdnote{} till \eighthnote{} are available for further notes

%\simplesumm{} % Simple summary

% Abstract (Do not insert blank lines, i.e. \\) 
\abstract{
Archived seismograms recorded in \nth{20} century present valuable source of information for monitoring earthquake activity and geophysical modelling. However, old data are only available as scanned paper-based images that should be digitised and converted from raster to vector format. Seismograms have special characteristics recorded by seismometer and encrypted in the images: traces, minute time gaps, timing and wave amplitudes. This information should be recognised and interpreted. Our objective was to automatically digitise historical seismograms recorded by Galitzine seismometer in 1954 in Uccle seismic station, Belgium. A dataset included 145 TIFF images. Software for digitising seismograms are limited and many have disadvantages. We used DigitSeis and performed a full workflow which included pattern recognition, classification, digitising, corrections and converting TIFFs to digital output. Selected images were damaged with identified defects (spots, overlaps of traces). Such files were corrected interactively. The generated contours of signals were presented as time series and converted into digital MATLAB format (.mat files) which preserved the ground motion information contained in analog seismograms. The DigitSeis approach demonstrated a powerful and reliable toolset for processing analog seismic signals. This study contributed to methodological development of signal processing for seismograms archived in Royal Observatory of Belgium.
}  

% Keywords
\keyword{seismology; Galitzine seismometer; horizontal component; analogue seismogram; digitising; earthquake recording; ground motions; historical seismograms; seismic waves} 

\begin{document}

%\begin{paracol}
%%%%%%%%%%%%%%%%%%%%%%%%%%%%%%%%%%%%%%%%%%

\section{Introduction}

\subsection{Background}
The seismicity of the Earth represent a physical phenomena resulting from the tectonic processes of energy accumulation and release \cite{GutenbergRichter}. This fundamental concept of the Earth's physics is reflected in the movements on the surface that have different intensity of the vibrations, caused by the fluctuations of energy \cite{Shearer2019}. In geodynamics, the seismicity is the result of self-organisation of the Earth, responding to the lithosphere movements \cite{UdiasBuforn2017}. Measuring seismic signals using seismometers enables to evaluate ground motions of the Earth’s, associated with the earthquakes of various magnitude or ambient seismic noise \cite{nakata_gualtieri_fichtner_2019}. Evaluating seismic ambient noise is effective in different contexts. Examples include imaging and monitoring the Earth’s interior, seismic mapping from global to regional or local scales, environmental studies, polar research on elastic deformations from icequakes, to mention a few.

The importance of seismic datasets can be illustrated by their practical application in the fields of tectonics, geology and the physics of the Earth. Specifically, the wide use of seismograms and their features \cite{Ammon2021} includes such fundamental domains as geophysics \cite{Kulhanek}, subsurface exploration, global and planetary tectonics \citep{Crampin,Jackson2002,Leroyetal2010}, risk assessment in civil engineering \citep{Phillips,Pileckietal2017, Kurzeja}, analysis of earthquakes and relative features (dynamic rupture, seismic tomography and wave forms) \cite{Ahorner,Steinmann,Amorese,Lomnitz,Yuliatmoko}, or volcanology \cite{Pezzo}. Recently, with the advances in studies of the ambient wavefield, or "noise"-based techniques, it has become possible to use the data recorded between the earthquakes to conduct monitoring \cite{Ratdomopurbo,Mordret} or imaging \cite{brenguier2014mapping,Lemenkova202012,Braszusetal2021,Lemenkova202016,Lietal2000} in various contexts. 

Oceans are a global source of seismic energy \cite{Ardhuinetal2011,AucanArdhuin2013,PeureuxArdhuin2016}. They mostly generate surface waves that travel large distances and can be exploited for the aforementioned studies. The microseisms are ubiquious on Earth, and in most cases, are the dominant component of the ground motions recorded at the surface. In turn, the study of microseisms allows studying the dynamic of the atmosphere-ocean-solid earth couplings \cite{Bromirski1999,Williams,Ardhuinetal2011}, and provide a unique insight into the evolution of oceanic climate \cite{Singh2021,Toyokunietal2015}.

Such analyses are largely based on the retrieval of information from digital seismograms \cite{Claerbout1992}. Uniformed, standardised, converted to digital format and accurately processed seismograms enable to explain the physical nature of the ground motion signals. However, the use of such information relies heavily on our access to large datasets of seismograms. 
The valorization of such datasets depends on  accurate and standardised recording systems using commonly accepted and approved instruments (seismometers and seismographs) with products that remain consistent through decades of operation \cite{UdiasBuforn2017}. 

\subsection{Historical development}
The history of seismology has undergone a long process of technical evolution of \ac{EO} data processing \cite{DeweyByerly1979,Agnew}. The instruments used for seismic data recording went through a constant development of technologies and state-of-the-art methods. The systematic, precise and operative observation of the seismicity of the Earth began in late \nth{19} century along with a progress in technology, development of new seismographs and general advances in geophysical methods \cite{SchweitzerLee2003}. Since 1880s seismologists started recording ground motion continuously and systematically using seismographs \cite{SchweitzerLee2003}. From that period and until 1960s, seismograms were recorded in analogue format using technical abilities of the narrow-band low-range instruments of those time \cite{Batlloetal2008}.

During the first half of the \nth{20} century, seismograms were recorded in analog format on  paper \cite[e.g., ][]{Wangetal2014} or as magnetic tapes \cite{Anetal2015}. Starting 1960s, seismology benefited from the the onset of the computer era that led to a transition to digital data which, with instrumental development, increased in quantity worldwide \cite{Engdahl2002}. For instance, an unprecedented initiative of global seismological network \ac{WWSSN} was created in 1960s to generate a collection of the high-quality big seismic datasets \citep{OliverMurphy,Ammon2010}. This initiative was possible due to the uniformly accepted technical equipment and standardised workflow for data calibration in each station. 

Currently, seismograms are recorded and stored digitally in high-performance data centres using advanced algorithms of signal processing \cite{Heiner}. However, the remaining historical datasets from pre-degital times form a valuable database containing unique information on seismicity of the Earth that could be used for practical scientific purposes \cite{Kanamorietal2010,Bungum,Kanamorietal2019,Lomnitz2004,Ambraseys1983}. 

\subsection{Motivation}
Old analog seismograms present a key source of information regarding the seismicity of the Earth's in the pre-digital age. Therefore, there is also a time sensitive component of digitising them before they are degraded and lost over time. The importance of the historical datasets for geophysical research has been widely discussed \cite{Vogtetal1985, TevesCosta1999,Bono2007,Okal} and the urge to ensure their preservation stressed \cite{Okal}. Using archived information it is also possible to analyse past seismicity using historical texts, news or testimonies \cite{Alexandre1990a,Alexandre1985,Alexandre1990b}, which can be used for geophysical reconstructions and climate modelling. The analysis of seismic velocity as time series analysis can be used for monitoring volcanoes. Moreover, retrospective analysis of the old seismograms enables to model earthquake cycle as a continuous process and to get insights into the physics of the earthquakes. Finally, interpretation of old seismograms provides information for the long-term prognosis: predictive modelling could be possible by statistical analysis of the datasets using information on earthquakes's locations and magnitude \cite{Kanamori1988}.

\subsection{Objective and goals}
The goal of our study was to perform an automated digitisation of seismic traces contained in historical seismograms. To achieve this goal, the objective was to apply the DigitSeis software for accurate semi-automated vectorising of seismograms. The aim was to convert an existing large historical archive of seismograms (over 140 scanned raster images) from TIFF format into a digital MATLAB output (generate .mat format files) using advanced technical functionality of machine processing as automatically as possible. As a test case, we aimed at revitalising the historical dataset of old seismograms recorded in Uccle station 1954 by Galitzine seismometer. Upon the workflow of scanning and archiving Galitzine data we used the DigitSeis instrument that was  tested for vectorising scanned seismograms. Our study contributes to the existing publications on seismic signal processing using archive seismograms from \ac{ROB}. We also try to identify the limits of the tool to later see how the gap can be bridged using ML tools for a more comprehensive result.

\subsection{Approaches}
The first attempts to digitise the analogue recordings of seismic events into numeric format started in 1960s. An \emph{ad hoc} algorithm of thinning the lines on scanned seismograms and converting the image into a 0-1 digital matrix was developed by \cite{Jarosch1975}. Another study used a device initially used for weather maps and adjusted for automatic converting of seismic signals into a digital form, and plotting them using a cathode-ray tube display \cite{AdamsAllen1961}. The progress of digitising signals accelerated in 1980s \citep{AspinallLatchman,CrouseMatuschka1983} and 1990s \citep{Scherbaum1996a,NorlundGriffiths1993,McGee1995,Harjesetal1993,Scherbaum1996b,Scherbaum1996c} along with the increase of computer-based algorithms and advanced geophysical software. The digitising approaches have become more and more automated since 2000s, along with a rapid development of the programming languages and IT tools \cite{Bietal2016,Ishiietal2015,Schenke,Wangetal2014}. 
 
Recently, the use of the automated digitisers of seismograms or accelerograms (recording the acceleration of the strong motion ground) instead of manual vectorisers has received much attention. This includes the automated recognition of noise \cite{CrouseMatuschka1983}, handling smooth and wigged traces by algorithms for tracing waves based on pixels context \cite{Bietal2016}, synchronisation of the time scale for the three motion components, handling scratches, distortions or line crossings on the image, adjusting rotated position \cite{Trifunacetal1999}. 

Digitising old seismograms is a not straightforward task. Often the problems are caused by the recording methods in old mechanical seismographs. For instance, high velocity of stylus that does not touch the paper, clipping of amplitudes, low-contrast dark background on a smoked paper, records damaged over time due to bad storage conditions of paper \cite{Baskoutasetal2000}, to mention a few. At the same time, digitising process should be as automated as possible, to avoid human-induced errors. 

As a response to the arising needs for automatic, quick and accurate digitising and modelling of old seismic data, various programs were developed. The advanced in pattern recognition are supported by rapid growth of the programming languages \cite{Esmaili,analytica2030006,Lemenkova201991,Lemenkova2019100,Lemenkov202103}, and specifically, Python \cite{Cox,Stoneback, Lemenkova201989}, improvements of algorithms of \ac{ML} and \ac{DL} for image segmentation and clustering \cite{Jenkins,signals1010005,DebeirDecaestecker,signals2030037,Debeir2008}, and signal processing \cite{Foucart,Debeiretal2019}. 

Software for digitising seismograms is based on pattern recognition. For example, the Teseo \cite{Pintoreetal2005} traces raster files using cubic Bézier curves. This vectoriser handles monochrome images by manual and automatic methods and neural networks. A software using C\# uses algorithm of colour scene recognition by Color Scene Filed Method \cite{Wangetal2016}. A program using MATLAB GUI used algorithm of inversion problem for a matrix of model parameters and observed seismic data \cite{XuXu}. Another example of seismic software, the SeisDig by Berkeley Seismological Laboratory presents an interactive MATLAB GUI based tool performing trace inspection, checking time conservation, recognising events and inputing metadata header \cite{BromirskiChuang2003}. An approach to seismic recording is presented in DigiSeis \cite{Khanetal2012} aimed at digitizing signals using a built-in PC sound card for signal processing. It visualises the two channel seismic data in time and frequency, and enables to perform a time series analysis. Extracting wave data from paper seismograms was attempted by \cite{Liuetal2001} with a seismogram digitization and DB management system using Delphi3 based on Pascal. These are some examples of the existing software for digitising seismograms. Besides that, there are also pilot seismic studies with computer vision, graphical plotting and mapping in geophysical domains and seismology \cite{Lemenkova202017,Lemenkova202032}.

%%%%%%%%%%%%%%%%%%%%%%%%%%%%%%%%%%%%%%%%%%
%\newpage
\section{Materials and Methods}

\subsection{Study Area}
The study was performed in the Université Libre de Bruxelles, École polytechnique de Bruxelles, Laboratory of Image Synthesis and Analysis (LISA), Belgium. We used archives obtained as a courtesy of \ac{ROB}, Department of Seismology \& Gravimetry, \ac{UCC}, Fig. \ref{fig:1}. The process is done within the framework of the SeismoStorm project.

\begin{figure}[H]
        	\centering
        		 \includegraphics[width=13.0cm]{fig_01.jpg}
        		 \caption{Location of Uccle station on a topographic map of Belgium. Map source: authors.}
     \label{fig:1}
\end{figure}

Established in 1899 by Eugène Lagrange, Uccle was the only seismic station in Belgium for the most of the \nth{20} century. Afterwards seismic monitoring gradually expanded to a larger seismic network system of Belgium \cite{Camelbeeck1985} with Uccle remaining its important centre. 

\subsection{Instrument}
The raw dataset obtained from the Galitzine seismometer were processed using DigitSeis vectoriser. The original seismograms have been received from \ac{UCC} using the Galitzine seismometer, Fig. \ref{fig:2}. It is the \nth{1} electromagnetic instrument developed by B.B. Galitzine in 1910 (Fig. \ref{fig:2}) \cite{Somville1922a,Somville1922b}, designed to record a single component (horizontal or vertical) for earthquake recording \cite{Wgg}. This type of seismographs is constructed using a following principle: a coil of wire is fixed to the pendulum which oscillates between the poles driven by the earthquake forces. The movement of coil between a strong magnet generate an electric induction current. The current is being transferred to a sensitive galvanometer to direct a beam of light onto photographic paper \cite{Neumann}. 

\begin{figure}[H]
      	\centering
         		\includegraphics[width=8.0cm]{fig_02.jpg}
     \caption{Instrument used for data capture in 1954: Horizontal Galitzine seismometer located in {UCC}. Image source: courtesy of \ac{ROB}. Photo: Rapha\"el S. M. De Plaen.} \label{fig:2}
\end{figure}

The structure of the Galitzine seismometer is presented in Fig. \ref{fig:2}: the pendulum (at the front), consists of a mass suspended by the two wires. It includes several parts: copper plate, fixed on the rod next to the coil; an electromagnetic seismometer with a coil attached to the arm of the pendulum oscillating in a magnetic field between the poles of a magnet (at the back) and generating an electric current by induction; Foucault current for damping; horizontal component (7kg); a tracing device with an optical recording system including a photographic paper and a galvanometer mirror with a mobile frame; and the speed settings (period of 12s). 

B.B. Galitzine improved the accuracy of the seismometer by reducing the sensitivity of the system through diverting a part of the current. Thus, special features of Galitzine seismometer include the following: (1) Damping the resistance to the galvanometer (2) Damping the seismometer pendulum through the cumulated effect of a magnetic damping device and the external resistance to the pendulum coil. Due to such effectiveness, such principle is used for measurements by seismometers of the Galitzine's type \cite{ChakrabartyTandon}.

\subsection{Workflow}
\begin{wrapfigure}{R}{6cm}
  \centering
    \includegraphics[width=6cm]{fig_03.png}
  \caption{Conceptual flowchart for the technical process. Plotting: PlantUML.}\label{fig:03}
\end{wrapfigure}

The general flowchart scheme for the process used in this study is schematically visualised in Fig. \ref{fig:03}. This study is based upon an integrating approach combining several programs for data processing. 

The data were collected using existing archived of \ac{ROB} using Galitzine seismometer in 1954. The data were recorded by the drum of the seismometer and stored in boxes in a storage space of the ROB. In 2021 the data were scanned using scanner in A0 format in a very-high resolution (600 dpi) and saved as TIFF files. The TIFF files were then imported to the DigitSeis where the program converted them into 8-bit native HSV format.  

The next steps include the processing of the DigitSeis that has several step itself (Fig. \ref{fig:05}). the most important of them include pre-processing (e.g. image loading, measuring time gaps, navigating selected segments), object classification (e.g. setup of image threshold, marking time parameters, small region analysis for detailed trace exploring), vectorisation, digitising and post-processing that includes converting the output file into the .mat format. 

Most of the traces were well recognised by the program in the process of digitising using embedded algorithms, while selected segments required interactive corrections. In this way the steps performed in DigitSeis were repeated iteratively for each seismogram. The results of DigitSeis data processing shown a relatively good recognition of lines and traces. The validation of data was performed for file export.

The resulting digitised seismic traces have a series of traces and time gaps identified by the program automatically. We converted the file into the MAT format and processed the next steps using Matplotlib library of Python. In Python we used Matplotlib library for post-processing steps and quality control. We analysed the repeatability of the segments broken by the DigitSeis and visualised vector traces received from the DigitSeis as separate plots. We also checked the time ticks automatically recognised by the program per minute using time gaps, and segments of the lines. 

%\newpage
\subsection{Data}
The dataset received from the Uccle station presents a large collection of scanned paper-based seismograms (selected examples are presented in Fig. \ref{fig:04}). 

\begin{figure}[H]
     \centering
     \begin{subfigure}[h]{0.7\textwidth}
      	\centering
         	\includegraphics[width=13.0cm]{fig_04a.png}
         	\caption{Empty records between the lines of seismic traces with enlarged fragment of seismogram. Here: UCC19540106Gal\_N\_0811.TIFF}
     \end{subfigure}
     
  \vspace{3mm}
     \begin{subfigure}[h]{0.7\textwidth}
         \centering
        		 \includegraphics[width=13.0cm]{fig_04b.png}
        		 \caption{Partially spotted image caused by storage, with enlarged fragment of seismogram. Here: UCC19540107Gal\_N\_0815.TIFF}
     \end{subfigure}
          
 \vspace{3mm}
     \begin{subfigure}[h]{0.7\textwidth}
         \centering
        		 \includegraphics[width=13.0cm]{fig_04c.png}
        		 \caption{Continuous noise dark background on the image with blurred traces which cause lack of contrast for automated image recognition. Here: UCC19540108Gal\_N\_0815.TIFF}
     \end{subfigure} 
     
  \vspace{3mm}
     \begin{subfigure}[h]{0.7\textwidth}
         \centering
        		 \includegraphics[width=13.0cm]{fig_04d.png}
        		 \caption{Overlapped traces of a very large event which cause a problem during vectorising with selected enlarged fragment of seismogram. Here: UCC19540112Gal\_E\_0750.TIFF}
     \end{subfigure}
     \caption{Examples of the raw data: old analog seismograms recorded by Uccle seismic station, Belgium, 1954, scanned and saved as TIFF files. Some images have visible distortions and defects.} \label{fig:04}
\end{figure}

The final images are stored as TIFF files which present a matrix of 2D arrays of pixels where each pixel is represented by a number of grayscale corresponding to that pixel \cite{Poynton}. The images have been stored at a resolution of 300 dpi, ca. 350-400 MB size each, in 256 B/W monochrome scales. The standard format of all the images ensured the integrity and compatibility of the whole dataset used for analysis. The 300 dpi was considered enough for ambient seismic noise analysis, even though some files were originally saved at 1200 dpi. The records from the Galitzine seismometer are monochrome but those of Wiecherts are stored in color.

The dataset included a collection of 145 images covering the period from 1 January 1954 to 12 March 1954. Some of the images are well preserved yet some have distortions and defects visible on the aged paper, Fig. \ref{fig:04}. The scanned images of old seismograms have different features and characteristics, which increase technical difficulties of processing such data: line breaks, distortions, defects, blurred lines, spots, etc. For instance, some data had defects due to the record imperfectness or storage cases. Oftentimes the traces are slanted (e.g. Fig. \ref{fig:06}, \ref{fig:10}, \ref{fig:11}) while some others appear more horizontal Fig. \ref{fig:07}, which is caused by the physical features of spiral recording by a rotating drum in the seismometer. Technical tasks concerned the quality of the recordings on such data with each individual seismogram and challenges to automatically recognise traces on the old scanned paper. 

\subsection{Software and Workflow}
In this study we used the DigitSeis, a MATLAB based software for processing analog scanned seismograms \cite{BogiatzisIshii2016} using a structured workflow, Fig \ref{fig:05}.

\begin{figure}[H]
        	\centering
        		 \includegraphics[width=13.0cm]{fig_05.jpg}
        		 \caption{Pipeline schematic diagram for workflow process used in this study}
     \label{fig:05}
\end{figure}

The DigitSeis is developed as a software instrument for digitising seismograms obtained by scanning the paper. It applies the extracts information from the scanned seismograms processed as an image matrix. DigitSeis converts scanned raster images of analog seismograms from original TIFF format into the convenient digital format using signal processing algorithms. The approach of DigitSeis is based on the principles of semi-automated algorithms of signal processing as time series. It applies minor human supervision, mostly for adjustment of functionality in selected processes and monitoring technical workflow (Fig. \ref{fig:06}). 

An embedded algorithm in the DigitSeis uses trace analyses which enables to track the lines and record their values using time gaps. We applied this method to identify traces on the scanned TIFF files of seismograms and processes them by automatic detection of lines and distortions (broken lines, spots, repetitive gaps on the series of lines, etc). The digitising was adjusted to distinguish the presence of lines, similar to the '0-1' matrix, using methods of colour recognition on the image. However, the directions of vector lines present another task that can only be solved using advanced analysis of the geometric trend of the trace by \ac{ML} tools. The thickness of the traces and time gaps on the data were considered for each 30-min line on the seismograms. 

\begin{figure}[H]
     \centering
     \begin{subfigure}[h]{0.7\textwidth}
      	\centering
         	\includegraphics[width=13.0cm]{fig_06a.jpg}
         	\caption{Original seismogram (UCC19540109Gal\_E\_0812.tif)}
     \end{subfigure}
     \vspace{3mm}   

     \begin{subfigure}[h]{0.7\textwidth}
         \centering
        		 \includegraphics[width=13.0cm]{fig_06b.jpg}
        		 \caption{Seismogram processed by DigitSeis}
     \end{subfigure}
     \caption{Seismogram recorded from Uccle station on 9 January 1954} \label{fig:06}
\end{figure}

The workflow of seismogram vectorisation by DigitSeis can be summarised in five major steps as follows: 1) Image preprocessing; 2) Identification of time marks and traces; 3) Image vectorising and classification; 4) Trace correction and control; 5) Timing and conversion (SAC).

\subsubsection{Image loading and preprocessing}
First, the images were loaded into the program as TIFF files, converted into the native 8-bit \ac{HSV} formal of DigitSeis where the intensity of monochrome images is either 0 (black) or 255 (white) in terms of intensity. Visually, the image contained white traces of measured seismic signals over a black background (Fig. \ref{fig:06}). In the selected images with limited data and large dark background the areas of interest were cropped to reduce the space without traces and to select only a part of image to test a smaller space area of interest (Fig. \ref{fig:07}). The resulted analysis was saved continuously throughout the workflow into a in the binary data MATLAB format as a .mat file.

\begin{figure}[H]
     \centering
     \begin{subfigure}[h]{0.7\textwidth}
      	\centering
         	\includegraphics[width=13.0cm]{fig_07a.png}
         	\caption{Original scanned seismogram (UCC19540116Gal\_E\_0820.tif)}
     \end{subfigure}
     
  \vspace{3mm}
     \begin{subfigure}[h]{0.7\textwidth}
         \centering
        		 \includegraphics[width=13.0cm]{fig_07b.jpg}
        		 \caption{Enlarged fragment of image (a) showing the time gaps indicating minute marks and zoomed segment separating the trace line between each other (tiny white gaps breaking traces)} 
     \end{subfigure}
     \caption{Analysing seismogram in paper-based format: traces and time gaps} \label{fig:07}
\end{figure}

\subsubsection{Identifying time gaps on seismograms}
Interpretation of seismic reflection signals involves the identification of arrival times recorded in the time markers or gaps. The exact determination of time and coordinates of the signals is essential in seismology. 

Therefore, we evaluated the 1-second time mark dimensions which show the breaks of the line along its course. The recording contained time marks at every minute. The whole line represents, on average, 30 min intervals of ground motion signals (Fig. \ref{fig:07}). On the original TIFFs of the seismograms, the time is marked according to the \ac{UTC} time and, most of the times, recorded at 30 mm/min. Marking time was performed by measuring width and offset of the gaps in the traces and recording the values in the menu bar. Time intervals on seismograms can be represented in two ways: time marks (small vertical dash lines) or time gaps (small breaks between the lines). In this study we only had time gaps. In the DigitSeis program time marks are indicated as positive numbers, while time gaps as negative numbers. Therefore, the time intervals were indicated as negative numbers, Fig. \ref{fig:08}. 

Identifying time gaps is a preparatory step before classification, as time marks are used to properly identify breaks between the signal traces. In out dataset, the seismograms had only gaps for time marks, without marks, so we only needed to identify traces as necessary objects. To receive accurate mean measurement values, we indicated time gaps from -22 to -20, so that the gaps of smaller value could be included, Fig. \ref{fig:09}. After time mark width was set up, the signals were prepared for classification. The program identified traces based on the analysis of areas of contrast in a monochrome colour scale and recognised the position of traces. We considered the curvature of the traces and measured the time gaps according to the path direction. As a result, vertical time gaps were identified for which the threshold number was set such that it was enough to include the individual traces with smaller (or less distinguishable) gaps. Once we identified time gaps the seismograms were ready for classification. 

\begin{figure}[H]
     \centering
     \begin{subfigure}[h]{0.7\textwidth}
      	\centering
         	\includegraphics[width=13.0cm]{fig_08a.jpg}
         	\caption{Identifying time marks on seismograms by measuring time gap between records. Here: example on file UCC19540119Gal\_N\_0825.tif}
     \end{subfigure}
     \vspace{3mm}   

     \begin{subfigure}[h]{0.7\textwidth}
         \centering
        		 \includegraphics[width=13.0cm]{fig_08b.jpg}
        		 \caption{Indicating time marks on seismograms as -22 and preparing image for classification}
     \end{subfigure}
     \caption{Classification of seismogram processed by DigitSeis.} \label{fig:08}
\end{figure}

\subsubsection{Classification}

The aim of the classification algorithm is to discriminate the traces from the time marks and noises, and to perform a complete object recognition, i.e., all objects are to be classified and assigned to the correct category. The image threshold for classification was set as 82 Fig. \ref{fig:09}, representing the minimum intensity value of a pixel to be included in classified objects. In our case, 82 was appropriate number for moderately clear images and thin lines. The Classification was performed using 'Classify Objects' tool of DigitSeis and the objects on the image were classified into 2 categories: 'traces' and 'noise' (in our case we did not have time marks, but time gaps instead). In some cases, noise objects were misclassified into the 'trace' category and vice versa. Therefore, their categories were adjusted accordingly using the 'Change Classification' function. Here traces are separate segments in a seismogram interrupted by the time gaps. 

\begin{wrapfigure}{r}{0.28\textwidth}
  \begin{center}
    \includegraphics[width=0.25\textwidth]{fig_09.jpg}
  \end{center}
  \caption{Threshold parameters for seismogram classification in DigitSeis}\label{fig:09}
\end{wrapfigure}

These time gaps mean minute intervals in traces, necessary to calculate timing. Noise include all irrelevant objects and annotations on the seismograms, e.g. spots, handwritten texts or notations, usually written on the edges of seismograms. Using 'Classification' tool we converted the image from raster TIFF format into vector binary format and classified scanned image into a set of objects with 2 main categories: 1) traces; 2) noise Fig. \ref{fig:10}. This dataset did not have time marks, but time gaps instead; otherwise we would have a third class of objects -- 'time marks'. The classification algorithm performed the partition of image into classes and assignment of pixels into each categories. 

\begin{figure}[H]
     \centering
     \begin{subfigure}[h]{0.7\textwidth}
      	\centering
         	\includegraphics[width=13.0cm]{fig_10a.jpg}
         	\caption{Results of the classified seismogram with shown identified object categories (fragment): traces (white) and noise (red, here: handwritten annotations)}
     \end{subfigure}
     \vspace{3mm}   

     \begin{subfigure}[h]{0.7\textwidth}
         \centering
        		 \includegraphics[width=13.0cm]{fig_10b.jpg}
        		 \caption{Small region analysis used for defining a smaller area of interest for closer examination of a border region of the seismogram}
     \end{subfigure}
     \caption{Results of the paper seismogram classification processed by DigitSeis} \label{fig:10}
\end{figure}

The algorithm of classification identified lines and traces, and discriminated them from noise and time gaps. The results of the classification include two types of objects: traces and noise, coloured by red or white, where red signifies noise (not used in the vectorisation) and white - traces (signals of the ground motion, i.e. the majority of objects on the seismogram), Fig. \ref{fig:10}. Combining separate segments of the traces was done using the crosshairs, Fig. \ref{fig:11} (b). The classification settings were adjusted using the Classify Objects function of DigitSeis, in order to make the automated processing smooth and accurate (Fig. \ref{fig:11}). In this way, the procedure of image processing generated a pixel matrix for the traces by scanning the whole image.  

\begin{figure}[H]
     \centering
     \begin{subfigure}[h]{0.7\textwidth}
      	\centering
         	\includegraphics[width=13.0cm]{fig_11a.jpg}
         	\caption{Results of the classified image with shown yellow segments of the identified trace (enlarged fragment). Here: example of file UCC19540109Gal\_E\_0812.tif}
     \end{subfigure}
     \vspace{3mm}   

     \begin{subfigure}[h]{0.7\textwidth}
         \centering
        		 \includegraphics[width=13.0cm]{fig_11b.jpg}
        		 \caption{Classified seismogram with traces saved in binary format 0-1. Here: example of file UCC19540109Gal\_E\_0812.tif (seismogram recorded on January 9, 1954.)}
     \end{subfigure}
     \caption{Results of the scanned seismogram classification processed by DigitSeis} \label{fig:11}
\end{figure}

Analog seismograms almost always include imperfectness on the old paper (e.g. handwritten annotations, gaps in data, damaged parts of the records), technical nuances caused by the seismometer instrument. Therefore, to ensure detailed selection of objects, a closer analysis of a selected parts of the seismogram was done using the 'Small Region Analysis' function, Fig. \ref{fig:10} (b). As a result, traces, time gaps, and noise were all classified more accurately as a prerequisite for the next step of digitising. 

%%%%%%%%%%%%%%%%%%%%%%%%%%%%%%%%%%%%%%%%%%
\section{Results}

\subsection{Vectorisation}
In order to make the digitising workflow more effective and accurate, the images were zoomed for enlarged analysis. The example of the processed image taken on January \nth{9}, 1954 (UCC19540109Gal\_E\_0812.tif) contained 45 traces with numeration from top to bottom, each horizontal line corresponds to the 30 minutes of measurements, Fig. \ref{fig:12}. Thus, the whole image shows seismic situation for 24 hours taken during January \nth{9}, 1954 (for the case of UCC19540109Gal\_E\_0812.tif).

\begin{figure}[H]
     \centering
     \begin{subfigure}[h]{0.7\textwidth}
      	\centering
         	\includegraphics[width=13.0cm]{fig_12a.jpg}
         	\caption{Digitised traces upon completion of the automated vectorisation. Some time gaps in the upper left part of the image were too small and not clearly visible for automatic discrimination between the trace and dark background. In these cases, gaps required manual correction to identify time intervals.}
     \end{subfigure}
     \vspace{3mm}   

     \begin{subfigure}[h]{0.7\textwidth}
         \centering
        		 \includegraphics[width=13.0cm]{fig_12b.jpg}
        		 \caption{Enlarged view of the automatically recognised digitised traces displayed by lines of various colours, zero-lines for each trace (cyan, dashed lines, numbered from top to bottom) and time gaps (vertical yellow dashes)}
     \end{subfigure} 
     \caption{Results of the paper-based seismogram vectorisation processed by DigitSeis} \label{fig:12}
\end{figure}

For vectorising lines, DigitSeis uses a mathematical algorithm of 'golden-section search' based on the nonlinear optimisation approach which finds an extremum (min/max) of a function inside a given interval \cite{Kiefer,Press}. Thus it lessens the amplitude of the first derivative of the combined trace according to the relative vertical shift between the seismic trace and time gaps \cite{BogiatzisIshii2016}. The vectorisation was obtained using the embedded algorithm of DigitSeis. Examples of vectorisation of selected seismogram by automatic approach of DigitSeis are shown in Fig. \ref{fig:12}. The Fig. \ref{fig:12} (a) shows the vectorisation for the whole seismogram, in which many segments were identified correctly, while others required interactive corrections (encompassed in yellow boxes). The Fig. \ref{fig:12} (b) visualises the enlarged fragment of the same seismogram with automatically identified traces with curvatures and time gaps. 

\subsection{Post-Processing: Correcting Traces}

The majority of traces and time gaps were well identified automatically by DigitSeis, except for selected lines that required manual corrections in an interactive regime. Some of such traces on the images were classified wrong, Fig. \ref{fig:13}. 

\begin{figure}[H]
     \centering
     \begin{subfigure}[h]{0.7\textwidth}
      	\centering
         	\includegraphics[width=13.0cm]{fig_13a.png}
         	\caption{Identified wrong vector direction of line crossing individual traces}
     \end{subfigure}
     \vspace{3mm}   
     
     \begin{subfigure}[h]{0.7\textwidth}
         \centering
        		 \includegraphics[width=13.0cm]{fig_13b.png}
        		 \caption{Detected misclassifications caused erroneous digitising. The gaps on the zero-lines (small yellow boxes) show the gaps that existed in the old paper in the original image itself.}
     \end{subfigure}
     \vspace{3mm}   
     \caption{Identified traces for selective correction and re-digitising using Correct Trace mode} \label{fig:13}
\end{figure}

Mostly they were located in the difficult segments of the image where the program could not recognise selected traces correctly using algorithms of machine vision. These difficulties were caused by the entangled path of lines that could not be discriminated automatically from noise or where the machine could be recognise the correct path of the line. This resulted in erroneous digitising and required manual corrections of the parts of images that contain gaps (yellow boxes), Fig. \ref{fig:13}. 

\begin{figure}[H]
     \centering
        \begin{subfigure}[h]{0.7\textwidth}
         \centering
        		 \includegraphics[width=13.0cm]{fig_14a.png}
        		 \caption{Overlap of line traces unrecognised during digitising: one segment of trace went steeply downwards and merged with another trace}
     \end{subfigure}
     \vspace{3mm}   

     \begin{subfigure}[h]{0.7\textwidth}
         \centering
        		 \includegraphics[width=13.0cm]{fig_14b.jpg}
        		 \caption{Enlarged view of the manually corrected entangled traces}
     \end{subfigure} 
     \caption{Correcting misclassified traces with wrong direction based on colour and geometric pixel's characteristics. The pixels on the monochrome image were assigning to categories using threshold values to 1 (white) and all other pixels to 0 (black) for automated detection of lines.} \label{fig:14}
\end{figure}

Correcting traces in a seismogram was performed using the 'Correct Trace' mode of DigitSeis. The identified traces were examined in an enlarged view (Fig. \ref{fig:14}) and the lines were manually adjusted. Misclassified traces were corrected semi-automatically and the image was re-digitised. In this phase, we identified the location and the problem of the misclassified seismic traces on an image and classifies them anew into the segments. 

We used the magnifying tool for enlarging a small portion of the trace in a reclassification window, Fig. \ref{fig:14} and \ref{fig:15}. First, we identified traces that required correction, and then corrected the line accordingly. Afterwards, the updated trace information was integrated to the whole line of trace in each case. Some segments were manually changed in the misclassified traces and the types of the objects were corrected in these cases (e.g. annotations on the edges of paper have been assigned a category of 'noise'). Visual inspection was performed upon the re-classification.

\begin{figure}[H]
     \centering
     \begin{subfigure}[h]{0.7\textwidth}
      	\centering
         	\includegraphics[width=13.0cm]{fig_15a.jpg}
         	\caption{Merging the trace initially broken into the three separate parts (three small yellow boxes)}
     \end{subfigure}
     \vspace{3mm}   

     \begin{subfigure}[h]{0.7\textwidth}
         \centering
        		 \includegraphics[width=13.0cm]{fig_15b.jpg}
        		 \caption{Reclassification of the selected segment and digitising the centroid of the trace line}
     \end{subfigure}
     \caption{Correcting trace for the selected segments} \label{fig:15}
\end{figure}

The vectorised traces were then highlighted by random colours and corresponded to the traces in the image, Fig. \ref{fig:16}. The vectorised curves of traces were optimised by the approximations with minimum distortion of the line by splines. The results of the processing of the classified image Fig. \ref{fig:17} (a) are presented as digitised vector image with distinct coloured traces and minute time gaps randomly coloured on the digitised image, Fig. \ref{fig:17} (b). Correcting the classification of traces was a necessary step before vectorising. 

In general, the automatic tracing of lines by DigitSeis performed well for the segments with low amplitude that do not cross each other. For the cases where the traces overlap and interest due to the seismic-phase arrivals, large amplitudes cross neighbouring traces. In this cases we corrected the truncated waveforms manually as explained above. We used special tools allowing to pick up the trace, correct its vertical extent, and separate it from neighbouring lines by manually removing the parts of the segment that were erroneously classified to the given trace. Using zooming we isolated segments of seismic lines and performed classification anew, followed by time identification, and vectorising.

\begin{figure}[H]
     \centering
     \begin{subfigure}[h]{0.7\textwidth}
      	\centering
         	\includegraphics[width=13.0cm]{fig_16a.jpg}
         	\caption{Corrected digitised traces upon completion of the automated vectorisation}
     \end{subfigure}
     \vspace{3mm}   

     \begin{subfigure}[h]{0.7\textwidth}
         \centering
        		 \includegraphics[width=13.0cm]{fig_16b.jpg}
        		 \caption{Corrected digitised traces with visible SDT lines (bold) after automated vectorisation}
     \end{subfigure}
     \caption{Corrected digitised traces upon completion of the automated vectorisation} \label{fig:16}
\end{figure}

Manual correction in vectorising was performed by tracing down a line in a selected segments where merged regions were separated manually and lines directions were corrected along the necessary fragments of trace. Selected polygons were assigned manually by  vectorising using modality in the 'Correct Trace' function, Fig. \ref{fig:15}.

\subsection{Timing and Converting}

After all the traces were digitised, the time was calculated using 'Calculate Timing' function, Fig. \ref{fig:18}. The time calculation properties were setup using a menu where the program asks for the default time and time increment (in the demonstrated case, January 9, 1954), Fig. Fig. \ref{fig:18} (a). Thus, we set up the absolute time for one record as a reference point which was then highlighted in red. For the setup of time marks we identified the time gaps that correspond to the beginning of the minute and provided the absolute time using a special menu, Fig. \ref{fig:18} (b). It included the time in HH:MM format and year-month-day period that exist in each seismogram. The time and date associated with the time mark were entered and the time marks were recalculated, Fig. \ref{fig:18} (c).

\begin{figure}[H]
     \centering
     \begin{subfigure}[h]{0.7\textwidth}
      	\centering
         	\includegraphics[width=13.0cm]{fig_17a.jpg}
         	\caption{Classified image recognised by DigitSeis into 'traces' and 'noise' (Here: UCC19540311Gal\_E\_0727.TIFF)}
     \end{subfigure}
     \vspace{3mm}   

     \begin{subfigure}[h]{0.7\textwidth}
         \centering
        		 \includegraphics[width=13.0cm]{fig_17b.jpg}
        		 \caption{Distinct coloured traces and minute time gaps are visible on the digitised image (Here: UCC19540311Gal\_E\_0727.mat)}
     \end{subfigure}
     \caption{Classified image (a) and digitised results (b). Here: UCC19540311Gal\_E\_0727 (11 March 1954)} \label{fig:17}
\end{figure}

The duration for the most of the lines was 30 minutes and the seismograms were recorded by the drum of the seismometer in one day. So we indicated the time between the gaps in seconds, and the program drawn the gaps (visible yellow vertical dashed), Fig. \ref{fig:18} (c) and in the enlarged view,  Fig. \ref{fig:19}. The time markers were visualised using time gaps classified on the previous step and measured automatically using a threshold.

Afterwards, the resulted data were exported into the MATLAB format, and saved as the .mat file. In MATLAB environment vector file was opened, checked and overlaid over the scanned TIFF image. The resulting files converted into MATLAB contained the vector lines as traces, that is included information that we received in DigitSeis. The converted file included the vectorised image with traces, its classification structure with segments and parts of the lines and time gap breaks. Therefore, the special value of DigitSeis consists in compatibility with MATLAB and possibility of further reuse the file in other tools, like Python. In this way, the vectorising workflow that we performed using DigitSeis is a continuous procedure can be smoothly integrated and exported to another software as a stepwise chain of separated processes.

\begin{figure}[H]
     \centering
     \begin{subfigure}[h]{0.7\textwidth}
      	\centering
         	\includegraphics[width=13.0cm]{fig_18a.jpg}
         	\caption{Time calculation process}
     \end{subfigure}
     \vspace{3mm}   

     \begin{subfigure}[h]{0.7\textwidth}
         \centering
        		 \includegraphics[width=13.0cm]{fig_18b.jpg}
        		 \caption{Timing setup using time display increment}
     \end{subfigure}
     \vspace{3mm} 
     
          \begin{subfigure}[h]{0.7\textwidth}
         \centering
        		 \includegraphics[width=13.0cm]{fig_18c.jpg}
        		 \caption{Time markers at 1-minute intervals on each 30-minute trace}
     \end{subfigure}
     \caption{Seismogram image with adjusted timing. Here: UCC19540311Gal\_E\_0727.mat} \label{fig:18}
\end{figure}

\begin{figure}[H]
     \centering
     \begin{subfigure}[H]{0.7\textwidth}
      	\centering
         	\includegraphics[width=13.0cm]{fig_19a.jpg}
         	\caption{Scanning of trace lines on the digitised image (Here: UCC19540311Gal\_E\_0727.mat)}
     \end{subfigure}
     \vspace{3mm}   

     \begin{subfigure}[H]{0.7\textwidth}
         \centering
        		 \includegraphics[width=13.0cm]{fig_19b.jpg}
        		 \caption{Separation of traces and time gaps by filtering (time gap setup as '-12' in this case)}
     \end{subfigure}
     \caption{Fragments of the digitised image (Here: UCC19540311Gal\_E\_0727.mat) (11 March 1954)} \label{fig:19}
\end{figure}

\subsection{Validating Results in Python}
Validating the results and checking quality control of the digitised traces have been performed using Python \cite{van1995python}. For export and graphical visualisation of the MAT files we used library Matplotlib \cite{Hunter}. For tuning the parameters of the digitising process in the next series of recording we performed the post-processing of the result using graphical visualisation and the analysis of the performed vectorisation by exporting the file of MAT format into Python, Fig. \ref{fig:20} and \ref{fig:21}. 

The length distribution of the segments imported is represented in Fig. \ref{fig:21}, the peak correspond to the typical frequent segment length $S_{typ}$ (approximately 1415 pixels in our setup for 59 seconds of recording).
We arbitrarily keep as valid the segment inside the range $[.8*S_{typ},1.1*S_{typ}]$. The starting position of these segments is represented as green dots in Fig. \ref{fig:20} (a). 
From the distribution of short invalid segments, as represented in Fig. \ref{fig:20} (a), we can identify the small segments corresponding to the typical hours marks length $H_{typ}$ (around 210 pixels here). The valid hours marks are considered in the range $[.9*H_{typ},1.1*H_{typ}]$. The respective graphs are presented for the seismogram UCC19540311Gal\_E\_0727\_280.mat taken on March 11, 1954, Fig. \ref{fig:21} (c) and (d).

The starting position of the hours segments is represented by the blue dots in Fig. \ref{fig:20}. 
Most of the hours marks are aligned, since the rotation speed of the main drum is 2 revolution per hours, some marks having the same length are also present in the border and noise segments.
By simply looking to the color dot distribution, it is easy to check the overall quality of the detected segments. 

\begin{figure}[H]
     \centering
     \begin{subfigure}[H]{0.7\textwidth}
      	\centering
         	\includegraphics[width=13.0cm]{fig_20a.png}
         	\caption{Quality control for time gaps: missed marks in unrecognised segments}
     \end{subfigure}
     \vspace{1mm}   

     \begin{subfigure}[H]{0.7\textwidth}
         \centering
        		 \includegraphics[width=13.0cm]{fig_20b.png}
        		 \caption{Correctly identified time gaps controlled by Python's Matplotlib}
     \end{subfigure}
     \caption{Controlling digitising results using Python (Matplotlib library)} \label{fig:20}
\end{figure}

The misclassified pixels were induced by the distortions of paper which caused the erroneous recognition of pixels in the edge of paper (red dots). Such pixels are classified as noise in the segments of trace in a way that it erroneously depicts the edges of seismic trace, as shown in Fig. \ref{fig:20}. The average length of segments divided automatically by DigitSeis, assessed and visualised by Matplotlib, Fig. \ref{fig:21}. 

\begin{figure}[H]
\captionsetup[subfigure]{justification=centering}
% first row subplot
% include first image
\begin{subfigure}{.38\textwidth}
      	\centering
         	\includegraphics[width=7.0cm]{fig_21a.png}
         \caption{}
\end{subfigure}
\hfill  
% include second image
\begin{subfigure}{.38\textwidth}
         \centering
        		 \includegraphics[width=7.0cm]{fig_21b.png}
        	\caption{}
\end{subfigure}

% second row subplot
\bigskip
% include first image
\begin{subfigure}{.38\textwidth}
      	\centering
         	\includegraphics[width=7.0cm]{fig_21c.png}
         \caption{}
\end{subfigure}
\hfill  
% include second image
     \begin{subfigure}{.38\textwidth}
         \centering
        		 \includegraphics[width=7.0cm]{fig_21d.png}
        		 \caption{}
     \end{subfigure}
     \caption{Validating the results of segments tracing in Python, Matplotlib for selected seismogram. (a), (c): Frequency of segment length: size of various segments in a trace. (b), (d): Analysis of statistics on segment length per hours of recording. (a), (b): UCC19540109Gal\_E\_0812.mat (c), (d): UCC19540311Gal\_E\_0727\_280.mat.} \label{fig:21}
\end{figure}

Because the program classifies seismic traces and short time gaps based on their geometrical characteristics and threshold, the quality of scanned paper may affect the automatic recognition and interpretation. Therefore, some segments were merged into one continuous line, while others were broken apart. 

Fig. \ref{fig:21} illustrate the analysis of the performance of DigitSeis in vectorising segments of a trace for the two examples (UCC19540109Gal\_E\_0812.mat and UCC19540311Gal\_E\_0727\_280.mat): the average length of the segment identified by the program as 'trace' and the frequency of the segments per hour.  
Seismic lines were separated by the DigitSeis into segments according to the traced lines (a 30-minute trace on each seismograms). The program searched for the correct path of the line and zero-lines using parameters. The separating of traces and noise from the time gaps and background on the paper was performed automatically. Therefore, screening and examining each pixel of the raster image was performed through the defined classifier thresholds. This resulted in variations of segments of the traces. 

For example, the parameters for image classification included image threshold 82\%, time mark width for gaps as -12 (in some cases up to -22, depending on the quality of scanned paper). The setup parameters included threshold and intervals of time gaps corresponding to the minute marks made by the drum of the seismometer. Some gaps were identified better while other were blurred, Fig. \ref{fig:20}. As a results, discrimination of the traces differed based on admissible properties using thresholds, as assessed by Python, Fig. \ref{fig:21}. 

%%%%%%%%%%%%%%%%%%%%%%%%%%%%%%%%%%%%%%%%%%
\section{Discussion}

Using archived scanned seismograms has many applications in geosciences. For example, study past seismicity of the Earth, climate modelling, exploring the atmosphere-ocean-solid earth interactions, and many more. Archives of old data containing seismograms recorded during the pre-satellite period is a crucial source of information for such tasks. To fully utilise the potential of historical seismograms, the abilities and advantages of modern digitising systems should be used. In our research we demonstrated the functionality of DigitSeis for classification and vectorising old seismograms from the historical dataset of \ac{ROB}. Currently the collection includes data for 1954 but it is planned to include other scanned images for more recent years, based on availability. 

To do this, we introduced a classification driven approach of DigitSeis for processing analog seismograms. Using the functionality of this program we presented a workflow where we showed that it performs well on datasets without time marks (we had time gaps instead). We extracted the features of objects, automatically detected on seismograms and derived the vector file containing classified objects (segments of line traces, noise and time gaps). The algorithms of DigitSeis processed well relatively large files (each scanned seismogram has a size of ca. 350-400 MB) and converted these data into .mat files. Some misclassified segments were presented in complex regions of seismograms with overlapped traces where the corrections required manual adjustments. Graphical representations of the encountered problems and complicated cases are illustrated in relevant sections. 

The process of classifying and digitising paper-based seismograms is based on computing features of the scanned objects (lines and other marks, normally classified as noise). The major principle of DigitSeis is creating the decision on classifying training lines (traces of seismograms) based on threshold test. For instance, we set up experimentally the time gaps from -22 to -12 to indicate the small distance breaking the minute marks, which was used by the machine to identifies all those gaps through out the image. Based on the setup parameters of threshold, the machine identified the objects on seismograms by a step-by-step analysis for the whole image.

This workflow arrived at classification of the image and digitising the objects (converting the output file into the .mat format). However, certain segments of the lines required manual tracing the paths and corrections, as demonstrated previously. The most complicated cases (overlapped traces) were studied separately to gain more insights in the visual characteristics of these segments and to discriminate the traces from noise signals. To do so, we zoomed selected enlarged areas on the seismogram and performed corrections manually using a crosshair. Specifically, we identified broken segments, crossed directions of the lines and overlay in traces. 

We presented the practical application of DigitSeis software for digitising scanned historical seismograms in a semi-automatic regime. We note that while using a good digitiser is a major task of any seismological research, there are still more algorithms that need to be adjusted for a fully automated system of classification and digitising of seismograms. Processing seismograms is an essential part of geophysical research having its specifics compared to simple images: minute time marks made by the drum, seismogram timing, recognition of line directions, smooth data conversion, etc. Our goal was to present here a practical application in signal processing by DigitSeis using dataset on historical seismograms with time gaps.

%%%%%%%%%%%%%%%%%%%%%%%%%%%%%%%%%%%%%%%%%%
\section{Conclusions}

As a central added value of this paper, we demonstrated that DigitSeis software is a very used tool for digitising signal seismic traces by applying the state-of-the-art software to a new corpus of data, based on digitised historical records from \ac{UCC} station, \ac{ROB}. The aim of this research was to digitise scanned images from the historical seismogram archives recorded by the Galitzine seismometer in 1954 in \ac{UCC}, Belgium. Currently the dataset covers year 1954, however, the data collection with be expanded and gradually cover more recent periods, until now. We presented the workflow of digitising seismograms by DigitSeis and illustrated performed practical steps. General performance of DigitSeis shown reliable results in vectorisation of analog scanned seismograms. DigitSeis performed a very impressive job for most of the clean traces. 

Detailed user experience presented in this study also illustrated that this software limits making the fully automatic analysis impossible. For example, when it comes to special cases often in historical seismograms (noises or stains, crossing lines), a manual assistance is required making fully automated approach impossible. In such cases, functionality and operations in vectorising were limited and required interactive corrections of recognised traces. The tools of DigitSeis still require improvements for processing of big data as streams. For the fully automated data processing and rapid building of big data corpus of digitised seismograms, the advanced solutions of \ac{ML} techniques are useful. As a suggestion for this, we propose a \ac{ML} solution of combining DigitSeis with Python. 

Automated digitising of large seismogram datasets will enable to re-use archives of historical seismograms as a valuable source of information in computational geophysics. The subsequent generation of the digitised corpus on seismograms will enable to restore seismic information contained in historical records from \ac{UCC} to analyse characteristics of traces and parameters of ground motions. Using scanned monochrome seismograms from the Galitzine seismometer (e.g. Fig. \ref{fig:04}), we employed a DigitSeis to vectorise the signals of seismic traces, recognised using defined thresholds and parameters of time gaps, and then converted the vector files to .mat format.  

A corpus of digitised old scanned seismograms can be effectively used in geophysical research for earthquake modelling as retrospective datasets. If scanned seismograms are digitised rapidly and accurately, they can be released in an accessible manner for further analysis. For this aim, a copious amounts of historical seismic data must be processed, classified, digitised and annotated for further integration and comparison of old and current seismic situation which would enable analysis of the seismicity of the Earth. Once tested, the workflow for digitising seismograms can be exported to larger datasets and cover seismic archives since 1910s up to now to help building a corpus of digitised seismograms. 

Earthquakes are among the most important types of the world's geological hazards significantly affecting environment, human's lives and property. The analysis of ground motions of the Earth is a background for prognosis of earthquakes, which requires the use of the advanced software for effective processing of seismic data using \ac{ML} approaches of signal recognition and interpretation. Analog historical seismograms are valuable source of information for retrospective modelling of the Earth's seismicity. 

To meet these challenges and make geophysical research more accurate and thorough, effective approaches in processing analog seismograms are needed urgently. Some of the existing approaches in vectorising seismograms are presented in the proprietary software and can only be achieved through restricted access. These methods have their advantages and shortcomings, and the use of these software is not at the level required for rapid and accurate processing of analog seismograms. Reasons for this are a lack of fully developed open algorithms, finely adjusted to digitise specifically seismic data, and the fact that the geophysical digitising software for signals processing may not be available to through open access. 

In this study we demonstrated the successful approach of automated digitising seismograms using DigitSeis and pointed at advantages and drawbacks. By employing its methods the intelligent automated analyses of seismograms was performed. It considered specific format and characteristics of seismic data (amplitude of traces, time gaps, 30-minute record intervals). DigitSeis presented a possibility for reuse the pool of historical seismic data The approach of DigitSeis implemented algorithms on semi-automated digitising of seismic traces and detecting time gaps based on the threshold parameters. 

Integration of workflow with \ac{ML} is possible due to another advantage of the DigitSeis, which consists in conversion of generated file to the .mat format for post-processing in MATLAB. This enable data integration and overlay of vectorised traces over the scanned images for data control. The advantages of DigitSeis in digitising seismograms enable recurrent and consistent information capture on seismic signal levels, repeatability of the recorded peaks, time gaps, period of data capture, intensity of signals and strength of Earth's ground motion. All this information can be used for seismicity analysis and relevant geophysical studies. 


%%%%%%%%%%%%%%%%%%%%%%%%%%%%%%%%%%%%%%%%%%
\vspace{6pt} 

%%%%%%%%%%%%%%%%%%%%%%%%%%%%%%%%%%%%%%%%%%
\authorcontributions{Conceptualization, O.D.; methodology, O.D. and P.L.; software, P.L.; validation, O.D.; formal analysis, O.D. and P.L.; investigation, O.D. and P.L.; resources, P.L.; data curation; writing---original draft preparation, P.L.; writing---review and editing, O.D. and P.L.; visualization, P.L.; supervision, O.D.; project administration at ULB, O.D.; funding acquisition, O.D. All authors have read and agreed to the published version of the manuscript.}

\funding{This research received funding from the Belgian Science Policy Office (BELSPO) for SeismoStorm project. We thank Rapha\"el S. M. De Plaen and Thomas Lecocq for providing the dataset.}

\dataavailability{Data supporting reported results are the courtesy of Royal Observatory of Belgium, Department of Seismology and Gravimetry (OD2), used in this study as available archived datasets of seismograms. Data are available in the Cytomine workspace}

\acknowledgments{The authors thank the anonymous reviewers for their comments and suggestions on the initial version of the manuscript.}

\conflictsofinterest{The authors declare no conflict of interest neither any personal circumstances or interest that may be perceived as inappropriately influencing the representation or interpretation of reported research results.} 

%%%%%%%%%%%%%%%%%%%%%%%%%%%%%%%%%%%%%%%%%%

\printacronyms[name=Abbreviations, include=abbrev, heading=section*]

%%%%%%%%%%%%%%%%%%%%%%%%%%%%%%%%%%%%%%%%%%

\reftitle{References}

%=====================================
% References
%=====================================

\externalbibliography{yes}
\bibliography{Seis.bib}
\end{paracol}
%\nocite{*}


%%%%%%%%%%%%%%%%%%%%%%%%%%%%%%%%%%%%%%%%%%
\end{document}

